\section{数据策略：规模之外还有质量、结构与学习阶段}

上一章给出模型怎样计算，本章追问模型究竟从什么分布中学习。数据策略不是“多抓网页”这么简单：source mix、质量过滤、顺序、重复率和阶段性配比都会改变梯度。课程先用人类学习对比提出问题，再用两个具体研究项目展示怎样把直觉转成实验。

\subsection{为什么数据是 backbone 而不是燃料桶}

前一章说明 block 怎样处理 token；现在必须把视线移到产生梯度的样本分布。把数据称为 backbone，是因为模型会把高频出现的事实、格式和策略内化为默认行为，而不是训练结束后再像数据库一样逐条查询。读这一节时先区分 corpus 的“数量”与“有效监督密度”：同样一个 token，如果重复、错误或与目标无关，它对能力的边际贡献可以完全不同。

对数据分布 $q(x)$，训练目标实际优化的是
\[
\mathbb{E}_{x\sim q}\big[\ell_{\theta}(x)\big].
\]
改变 $q$ 就改变模型被反复奖励的知识、风格、语言与推理模式。数据量增加只扩大样本数；若噪声、重复或域偏差同步增加，额外 token 可能产生递减收益甚至覆盖已有能力。


\lecturefigure{slide-024.jpg}{数据对 AI 的影响：质量、计算、规模与应用之间的耦合}{CS25 V5 official deck, p.~24}


\lecturefigure{slide-025.jpg}{更聪明的数据策略：质量优先、阶段配比与主动选择}{CS25 V5 official deck, p.~25}

\begin{knowledgebox}{读图：数据影响链应该怎么看}
先沿 slide 24 的因果链从 data 走到 model behavior，再读 slide 25 的策略列表。第一步比较 quality 与 diversity：前者减少噪声，后者避免能力收窄。第二步看 selection 是训练前过滤、训练中 sampling，还是训练后 active collection；三个阶段的成本不同。第三步看 metric 是否与目标能力一致。图支持“数据决策影响模型”，却不能从一张流程图推出某个过滤器在所有 scale 上都有效。
\end{knowledgebox}

\begin{importantbox}{数据策略的五个旋钮}
\begin{enumerate}
  \item \textbf{Quality}：样本是否正确、连贯、信息密度高；
  \item \textbf{Diversity}：语言、领域、任务与表达是否覆盖长尾；
  \item \textbf{Structure}：对话、故事、代码、证明等序列形式；
  \item \textbf{Scale}：token 数与重复次数；
  \item \textbf{Schedule}：不同数据在训练早期、后期出现的比例。
\end{enumerate}
同一个 corpus 在不同 schedule 下可能得到不同模型。
\end{importantbox}

\subsection{人类与 LLM：差异是研究假设，不是类脑证明}

课程用 baby/ChatGPT 与 brain/neural network 对照提出一组开放问题。人类持续学习、带目标行动、接收多模态感觉并在社会环境中获得反馈；经典 LLM 主要在一次大规模离线预训练中优化 next-token objective。这个对比的价值是生成实验轴，而不是宣称生物机制已被理解。


\lecturefigure{slide-026.jpg}{Baby 与 ChatGPT、brain 与 neural network：两类学习系统的对照}{CS25 V5 official deck, p.~26}


\lecturefigure{slide-027.jpg}{人类连续学习、目标、多模态与结构化归纳和 LLM 单阶段训练的差异}{CS25 V5 official deck, p.~27}

\teachervoice{00:12:38--00:14:46，Steven 把 continuous、goal-directed、multimodal 与 hierarchical learning 列为人类可能更高效的原因，同时反复使用 ``potential'' 与 ``might''。讲义把这些保留为 hypotheses，不写成神经科学定论。}

\begin{warningbox}{数据效率比较的三个陷阱}
人类接收的监督不只有可计数字符；视觉、动作、奖励、社会互动与长期环境反馈都难以折算成 token。模型和人类的目标函数也不同。最后，``会说话'' 的能力边界并不等同。因此“儿童 token 少得多”能提出问题，却不能单独证明某种数据源就是充分解释。
\end{warningbox}


\lecturefigure{slide-028.jpg}{儿童故事与对话：高结构、强互动、面向学习者的数据}{CS25 V5 official deck, p.~28}


\lecturefigure{slide-029.jpg}{互联网到 ChatGPT：海量但异质的网页语料}{CS25 V5 official deck, p.~29}

\subsection{为什么研究小模型与小数据}

前面的 human/LLM 对照给出很多可能因素，但若直接用最大模型验证，预算只够做一两个点，任何结论都会与超参和随机性纠缠。小模型的科学价值在于把一个大问题分解成密集曲线：重复训练、替换数据、观察收敛和错误类型。它也是 deployment 目标，因为本地模型关系到隐私、延迟与可访问性；不过这两个价值要分开，不能因小模型易研究就假设其 scaling 排序必然等于大模型。

小规模实验降低复现成本，使研究者能运行多 seed、做更密集 ablation，并检查训练动态；它也支持 on-device、隐私和开放研究。但小模型结果不能直接外推到 trillion-token regime，必须区分“发现机制候选”与“证明大规模收益”。


\lecturefigure{slide-030.jpg}{研究小模型与小数据的效率、可解释性、开放性和科学价值}{CS25 V5 official deck, p.~30}

\begin{knowledgebox}{读图：小模型研究有两条价值链}
一条是\textbf{scientific proxy}：低成本多 seed、密集 ablation、观察 learning curve；另一条是\textbf{deployment target}：on-device、低延迟、隐私和可访问性。读 slide 30 时不要把两条链混合。一个方法在 100M 模型上机制清晰，不等于它在 70B 上排序不变；一个 3B 模型能本地运行，也不等于它适合回答高风险问题。正确结论是小规模提供更好的实验分辨率，并为部署建立独立 Pareto frontier。
\end{knowledgebox}

\teachervoice{00:13:55--00:14:46，老师把小模型价值连接到本地运行、可控性、开源可及性与认知研究。实践经验：小实验最有价值的产物是可复核趋势和失败案例，不是廉价地模拟一次大训练。}

\subsection{本章小结}

数据策略同时包含质量、结构、多样性、规模与 schedule。人类学习对比给出可测试维度，小模型则提供低成本实验场。后面两项研究分别回答：儿童导向语料是否更有效，以及高质量数据是否应集中在训练后期。

